\section{Convergence proofs}
\label[appendix]{app:convergence-proofs}

\subsection{Proof of the modified energy law}
\label[appendix]{app:energy-law-proof}

\begin{proof}[Proof of \cref{thm:monotone-mobility-dissipation}]
Take the inner product of \cref{eq:abstract-momentum-balance} with
\(\Delta\theta^{n+1}=M^{n+1}p^{n+1}\).  Weighted polarization gives
\begin{align}
&\frac{(p^{n+1})^{\mathsf T}M^{n+1}
  (p^{n+1}-\beta p^n)}{\eta}
\nonumber\\
&\quad=\frac{1}{2\eta}\left(
  \norm{p^{n+1}}_{M^{n+1}}^2
  -\beta^2\norm{p^n}_{M^{n+1}}^2
  +\norm{p^{n+1}-\beta p^n}_{M^{n+1}}^2\right).
  \label{eq:weighted-polarization}
\end{align}
Insert and subtract \(\norm{p^n}_{M^n}^2/(2\eta)\) to obtain
\begin{align*}
&\frac{\norm{p^{n+1}}_{M^{n+1}}^2
  -\beta^2\norm{p^n}_{M^{n+1}}^2}{2\eta}\\
&\quad=
\frac{\norm{p^{n+1}}_{M^{n+1}}^2-\norm{p^n}_{M^n}^2}{2\eta}
+\frac{\norm{p^n}_{M^n-M^{n+1}}^2}{2\eta}
+\frac{1-\beta^2}{2\eta}\norm{p^n}_{M^{n+1}}^2.
\end{align*}
Substitute the scalar and pullback identities in
\cref{eq:common-identities}.  Collecting the scalar and kinetic energy
changes yields \cref{eq:monotone-mobility-dissipation}.  Each remaining
term is nonnegative because \(S_m^n\succeq0\),
\(M^n-M^{n+1}\succeq0\), and \(0\leq\beta<1\).
\end{proof}

\subsection{Proof of the local stability criterion}
\label[appendix]{app:schur-proof}

\begin{proof}[Proof of \cref{thm:local-schur-stability}]
Write \(\widehat H_*=M^{1/2}H_*M^{1/2}\) and
\(\widehat B_*=M^{1/2}B_*M^{1/2}\). For the linearized variables
\(x^n=M^{-1/2}(\theta^n-\theta_*)\) and \(w^n=M^{1/2}p^n\),
\begin{equation}
  (I_d+\eta\widehat B_*)w^{n+1}
  =\beta w^n-\eta\kappa_*\widehat H_*x^n,
  \qquad x^{n+1}=x^n+w^{n+1}.
  \label{eq:linearized-momentum-map}
\end{equation}
Derivatives of the inverse coefficient multiply a zero right-hand side,
and variations of \(q/Q\) multiply the zero gradient. The invertible change
of state from \((x^n,w^n)\) to \((x^n,x^n-w^n)\) gives the matrix polynomial
\begin{equation}
\begin{aligned}
  \mathcal P(z)
  ={}&z^2(I_d+\eta\widehat B_*)\\
  &-z\bigl((1+\beta)I_d+\eta\widehat B_*
                     -\eta\kappa_*\widehat H_*\bigr)+\beta I_d.
\end{aligned}
\label{eq:local-characteristic-polynomial}
\end{equation}
Its determinant vanishes exactly at the eigenvalues of the state map,
including when \(\beta=0\).

If \(\mathcal P(z)\xi=0\) with \(\xi\in\mathbb C^d\setminus\{0\}\),
multiplication by \(\xi^*\) gives the scalar quadratic
\begin{equation}
  (1+\eta b)z^2
  -(1+\beta+\eta b-\eta\kappa_*h)z+\beta=0,
  \label{eq:local-rayleigh-polynomial}
\end{equation}
where \(b=\xi^*\widehat B_*\xi/(\xi^*\xi)\geq0\) and
\(h=\xi^*\widehat H_*\xi/(\xi^*\xi)>0\) are real. The leading
coefficient exceeds \(\beta\). The values of this polynomial at \(1\) and
\(-1\) are \(\eta\kappa_*h\) and
\(2(1+\beta)+2\eta b-\eta\kappa_*h\), respectively, both positive under
\cref{eq:local-stability-criterion}. Nonreal roots have squared modulus
\(\beta/(1+\eta b)<1\). For real roots, the same root product and the
positive values at \(1\) and \(-1\) exclude roots outside \((-1,1)\).
This proves sufficiency.

Conversely,
\[
  \mathcal P(-1)=2(1+\beta)I_d+2\eta\widehat B_*
                                 -\eta\kappa_*\widehat H_*.
\]
If this matrix is singular, \(-1\) is an eigenvalue. If it has a negative
eigenvalue, \(\mathcal P(-t)\) becomes positive definite for sufficiently
large \(t\), since its leading coefficient is positive definite.
Continuity of its smallest eigenvalue gives a singular matrix
\(\mathcal P(-t)\) at some real \(t>1\), and hence a state eigenvalue \(-t<-1\).
This proves necessity.
\end{proof}

\subsection{Comparison with the appended recurrence}
\label[appendix]{app:appended-comparison}

On a common invariant one-dimensional mode of \(\widehat H_*\) and
\(\widehat B_*\), let \(h>0\) and \(b\geq0\) be their curvatures.
If \(b=\kappa_*h\), the characteristic equation reduces to
\[
  (1+\eta b)z^2-(1+\beta)z+\beta=0.
\]
For \(\beta>0\) and \(\eta b>(1-\beta)^2/(4\beta)\), its roots have modulus
\(\sqrt{\beta/(1+\eta b)}\), which tends to zero as \(\eta b\to\infty\).
For \(\beta=0\), the roots are \(0\) and \(1/(1+\eta b)\).
The frozen appended recurrence has root product \(\beta\), so its spectral
radius is at least \(\sqrt\beta\). On modes with large \(\eta b\) the
momentum update therefore has the smaller spectral radius; this is a
stiff-limit statement, not a uniform rate comparison. Its stability
condition, however, is
\[
  \eta\kappa_*H_*
  \prec2(1+\beta)(M^{-1}+\eta B_*),
\]
which admits at least the interval of the momentum update for the same
pullback. It follows by diagonalizing the positive definite generalized
Hessian pair \((H_*,M^{-1}+\eta B_*)\). The interval enlargement in
\cref{cor:local-step-interval} compares component and aggregate pullbacks
within the momentum update.

\subsection{Local convergence with scalar relaxation}
\label[appendix]{app:local-convergence-proof}

\begin{proof}[Proof of \cref{thm:local-relaxed-convergence}]
Set \(z=(\theta-\theta_*,p)\) and \(Q_*=Q(\theta_*)\). Let
\(T(q,z)\) be the parameter--momentum update, without applying scalar
relaxation. This map is continuously differentiable near \((Q_*,0)\), and
\(T(q,0)=0\) for every nearby \(q\). Its derivative
\(L=D_zT(Q_*,0)\) is Schur stable by
\cref{thm:local-schur-stability}. The convergent series
\[
  P=\sum_{j=0}^{\infty}(L^{\mathsf T})^jL^j
\]
defines a positive definite matrix with
\(L^{\mathsf T}PL=P-I\). Hence \(\norm{L}_P<1\).
Continuity of \(D_zT\) and integration along the segment from \(0\) to
\(z\) give a cylinder
\[
  \norm{z}_P\leq a_z,\qquad |q-Q_*|\leq a_q<Q_*/2
\]
and \(r\in(0,1)\) such that
\begin{equation}
  \norm{T(q,z)}_P\leq r\norm{z}_P.
  \label{eq:local-uniform-contraction}
\end{equation}

The scalar must remain in this cylinder. Initialization and the acceptance
rule give the invariant \(0\leq q^n\leq Q(\theta^n)\). Within the cylinder,
uniformly in \(\rho\in[0,1]\), the increment of
\cref{eq:abstract-momentum-step} and the scalar increment \(d\) of
\cref{eq:common-d} satisfy
\[
  \norm{\Delta}+\norm{p^+}=O(\norm{z}_P),\qquad
  d=O(\norm{z}_P^2),\qquad
  0\leq\mathcal D=O(\norm{z}_P^2).
\]
Write the two branches of \cref{eq:generic-relaxation} as
\(q_{\mathrm{bud}}^2=(q+d)^2+\rho\mathcal D\) and
\(q_{\mathrm{cap}}^2=Q(\theta+\Delta)^2\). Then
\[
  |q_{\mathrm{bud}}^2-q^2|\leq c_4\norm{z}_P^2.
\]
Since \(q\leq Q(\theta)\) and \(\nabla F(\theta_*)=0\),
\[
  q_{\mathrm{cap}}^2-q^2\geq F(\theta+\Delta)-F(\theta)
  \geq-c_5\norm{z}_P^2.
\]
Both branches are bounded below by \(q^2-c_6\norm{z}_P^2\) with
\(c_6=\max\{c_4,c_5\}\), while the budget branch bounds their minimum
above. Thus
\[
  |(q^+)^2-q^2|\leq c_6\norm{z}_P^2.
\]
Dividing by \(q^++q\geq q\geq Q_*/2\) gives
\begin{equation}
  |q^+-q|\leq c_2\norm{z}_P^2.
  \label{eq:local-scalar-drift}
\end{equation}
This estimate covers either selected branch and branch ties.

Because \(\nabla Q(\theta_*)=0\),
\(|q^0-Q_*|\leq c_3\norm{z^0}_P^2\). Choose the initial neighborhood
so that \(\norm{z^0}_P<a_z\) and
\[
  \left(c_3+\frac{c_2}{1-r^2}\right)\norm{z^0}_P^2<a_q.
\]
Induction using \cref{eq:local-uniform-contraction,eq:local-scalar-drift}
then gives
\[
  \norm{z^n}_P\leq r^n\norm{z^0}_P,
  \qquad
  |q^n-Q_*|
  \leq\left(c_3+\frac{c_2}{1-r^2}\right)\norm{z^0}_P^2<a_q.
\]
The iterates cannot leave the cylinder, closing the induction. Scalar
increments are summable, so \(q^n\to q_\infty\geq Q_*/2\). The invariant
\(q^n\leq Q(\theta^n)\) gives \(q_\infty\leq Q_*\), and the drift
bound gives \cref{eq:local-scalar-limit}. Equivalence of norms yields
\cref{eq:local-geometric-convergence}.
\end{proof}


\subsection{Global stationarity estimates}
\label[appendix]{app:global-stationarity-proofs}

With \(L_B=\sup_n\norm{B_\alpha^n}\), the constants of
\cref{cor:best-iterate-stationarity} are
\begin{equation}
\begin{aligned}
  C_{\mathrm{mom}}
  &=\frac{4\eta}{\underline\kappa^2\mu_{\min}}
  \left[
    \left(\eta^{-1}+\mu_{\max}L_B\right)^2
    +\frac{\beta^2\mu_{\max}^2L_B^2}{1-\beta^2}
  \right],\\
  C_{\mathrm{dir}}
  &=\frac{\eta\left(\eta^{-1}+L_B\right)^2}{\underline\kappa^2}.
\end{aligned}
  \label{eq:stationarity-constants-families}
\end{equation}
One sets \(K=C_{\mathrm{mom}}\) under the assumptions of
\cref{thm:conditional-stationarity} and \(K=C_{\mathrm{dir}}\) under
those of \cref{cor:direct-stationarity}.

\begin{proof}[Proof of \cref{thm:conditional-stationarity}]
Sum \cref{eq:relaxed-energy-bound} and use
\cref{eq:strict-relaxation-budget} and \(\Hcal^n\geq0\) to obtain
\(\sum_n\mathcal D^n<\infty\).  The damping and inertial residual terms in
\cref{eq:monotone-mobility-dissipation}, together with the uniform lower
mobility bound, imply
\(
p^n\to0
\)
and
\(
p^{n+1}-\beta p^n\to0
\).
The uniform upper bound gives
\(
\Delta\theta^{n+1}=M^{n+1}p^{n+1}\to0
\).
Using these limits and the boundedness of \(B_\alpha^n\) in
\cref{eq:abstract-momentum-balance} yields
\(
(q^n/Q^n)g^n\to0
\), and the ratio lower bound gives \(g^n\to0\). Since \(F\in C^1\), every
accumulation point is stationary.
\end{proof}

\begin{proof}[Proof of \cref{cor:direct-stationarity}]
The direct update law \cref{eq:base-dissipation} and summability give
\(\Delta\theta^{n+1}\to0\).  The state
equation \cref{eq:base-pbsav-step}, the bounded pullback, and the ratio lower
bound then imply \(g^n\to0\).  The accumulation-point claim follows by
continuity of \(\nabla F\).
\end{proof}

\begin{proof}[Proof of \cref{cor:best-iterate-stationarity}]
For the momentum update, set \(e^n=p^{n+1}-\beta p^n\).  The state
equation, \(p^{n+1}=e^n+\beta p^n\), and the residual and damping terms in
\cref{eq:monotone-mobility-dissipation} give
\begin{align}
  \frac{q^n}{Q^n}g^n
  &=-\frac{e^n}{\eta}
  -B_\alpha^nM^{n+1}(e^n+\beta p^n),\nonumber\\
  \norm{g^n}
  &\leq\frac1{\underline\kappa}
  \left[
    \left(\eta^{-1}+\mu_{\max}L_B\right)\norm{e^n}
    +\beta\mu_{\max}L_B\norm{p^n}
  \right],\nonumber\\
  \mathcal D^n
  &\geq\frac{\mu_{\min}}{2\eta}
  \left(\norm{e^n}^2+(1-\beta^2)\norm{p^n}^2\right).
  \label{eq:gradient-dissipation-bounds}
\end{align}
The inequality \((a+b)^2\leq2a^2+2b^2\) gives
\(
\norm{g^n}^2\leq C_{\mathrm{mom}}\mathcal D^n
\).

For the direct update, \cref{eq:base-pbsav-step} gives
\begin{equation}
  \norm{g^n}
  \leq\frac{\eta^{-1}+L_B}{\underline\kappa}
  \norm{\Delta\theta^{n+1}},
  \qquad
  \mathcal D^n\geq
  \frac{\norm{\Delta\theta^{n+1}}^2}{\eta},
\end{equation}
and therefore
\(
\norm{g^n}^2\leq C_{\mathrm{dir}}\mathcal D^n
\).
Finally, summing \cref{eq:relaxed-energy-bound} and using
\cref{eq:strict-relaxation-budget} gives
\(
(1-\bar\rho)\sum_{n=0}^{N-1}\mathcal D^n\leq\Hcal^0
\).
This proves both estimates in \cref{eq:best-iterate-stationarity}; the
infinite-series claim follows by letting \(N\to\infty\).
\end{proof}
